
\documentclass[12pt,a4paper]{amsart}

\usepackage{amsmath,amssymb,amsthm}
\usepackage{mathtools}
\usepackage[colorlinks=true,linkcolor=blue!60!black,citecolor=green!50!black,
            urlcolor=blue!70!black]{hyperref}
\usepackage[shortlabels]{enumitem}
\usepackage{booktabs}

%% ---- Theorem environments ----------------------------------------
\newtheorem{theorem}{Theorem}[section]
\newtheorem{lemma}[theorem]{Lemma}
\newtheorem{proposition}[theorem]{Proposition}
\newtheorem{corollary}[theorem]{Corollary}
\theoremstyle{definition}
\newtheorem{definition}[theorem]{Definition}
\newtheorem{remark}[theorem]{Remark}
\newtheorem{problem}{Problem}[section]
\newtheorem{conjecture}[theorem]{Conjecture}

%% ---- Notation ---------------------------------------------------
\newcommand{\Kop}{\mathcal{K}_c}
\newcommand{\PN}{P_{N,c}}
\newcommand{\Pedge}{\Pi_{\mathrm{edge}}}
\newcommand{\Ev}{\mathrm{Ev}_N}
\newcommand{\AN}{A_{N,c}}
\newcommand{\Dedge}{D_{\mathrm{Edge}}}
\newcommand{\lmin}{\lambda_{\min}(N,c)}
\newcommand{\ip}[2]{\langle #1,\, #2\rangle}
\newcommand{\norm}[1]{\|#1\|}
\newcommand{\Eout}{\mathcal{E}_{\mathrm{out}}}
\newcommand{\PW}{\mathrm{PW}_\omega}
% IIKS vectors (from paper13cc)
\newcommand{\Fvec}{\vec{F}_c}
\newcommand{\Gvec}{\vec{G}_c}

%%=================================================================
\title[Edge Coercivity and the Airy Boundary Operator]
      {Edge Coercivity and the Airy Boundary Operator
       for Prolate Gram Operators:\\
       Uniform Spectral Separation at the PSWF Edge\\[4pt]
       \large (Paper~XXII in the Prolate--Weil Program)}

\author{Sebastian Schmalnauer}
\date{\today}

\begin{document}
\maketitle

\begin{abstract}
Paper~XXI established, via Lemma~$\Omega$ and Theorem~B,
that the boundary compression operator
$\Dedge := \Pedge(\AN - \Kop)\Pedge$
has at least one positive eigenvalue (existential edge obstruction).
The present paper addresses the uniform version:
\[
  \inf\,\sigma(\Dedge) \ge c_2(\delta,c) > 0.
\]
The proof strategy is organized as a \emph{transfer chain}
along the spectral axis Bulk $\to$ Edge $\to$ Airy.
The transfer rests on three structural pillars:
(i) an exact integrable (IIKS) representation of the
prolate concentration operator, which provides the only
RHP-level access to the edge scaling limit;
(ii) a two-level obstruction principle separating
structural (nonlinear phase consistency) from geometric
(Widom-to-Airy kernel transfer) components;
(iii) a Gram-coercivity reduction showing that
edge-frame stability implies a spectral-gap lower bound.
Conditionally on the spectral gap of the Airy operator
$D^{\mathrm{Airy}}$, we identify $\Dedge$ as a
\emph{Rand-Weyl operator} in the Airy scaling limit
and prove $\inf\,\sigma(\Dedge) \ge c_2^{\mathrm{Airy}}/2 > 0$.
\end{abstract}

\tableofcontents

%%=================================================================
\section{Position in the Program}\label{sec:position}
%%=================================================================

The Prolate--Weil program (Papers~I--XXI) arrived at the
following canonical endpoint (see \texttt{HEBELSTELLE.md}):

\begin{center}\small
\begin{tabular}{lll}
\toprule
Regime & Spectral statement & Status \\
\midrule
Bulk & $\sigma(\AN) \approx \sigma(\Kop)$ & Theorem~A (XXI) \\
Edge (existential) & $\exists\,\lambda > 0 \in \sigma(\Dedge)$ & Theorem~B (XXI) \\
\textbf{Edge (uniform)} & $\mathbf{\inf\,\sigma(\Dedge) \ge c_2 > 0}$ & \textbf{This paper} \\
\bottomrule
\end{tabular}
\end{center}

The gap between Theorem~B and the uniform statement is exactly:
\[
  \boxed{\exists \;\ne\; \forall \text{ on Edge}}
\]
not trivially but as: local instability vs.\ global geometric separation.

\subsection{The single open question}

From \texttt{HEBELSTELLE.md}, everything in the program is closed
except:
\[
  \text{Is } \Ev|_{\mathrm{Edge}} \text{ a uniform quasi-isometry?}
\]
or equivalently:
\[
  \boxed{\inf\,\sigma\bigl(\Pedge(\AN - \Kop)\Pedge\bigr) > 0}
\]

\subsection{Transfer structure of this paper}

The proof is organized as a chain along the spectral axis:
\[
  \underbrace{\text{Bulk}}_{\text{Thm.~A, XXI}}
  \;\longrightarrow\;
  \underbrace{\text{Edge}}_{\Dedge \ne 0,\;\text{Thm.~B, XXI}}
  \;\longrightarrow\;
  \underbrace{\text{Airy limit}}_{D^{\mathrm{Airy}} \gg 0,\;\text{this paper}}
  \;\longrightarrow\;
  \underbrace{\text{Uniform coercivity}}_{\inf\,\sigma(\Dedge) > 0}
\]
The three new structural inputs that make this transfer possible
are developed in \S\S\ref{sec:iiks}--\ref{sec:gram_gap}:
\begin{enumerate}[(I)]
  \item IIKS integrable representation of $\Kop$
  (\S\ref{sec:iiks}): provides the RHP-level access
  to the edge scaling limit that the Fourier representation alone
  does not afford.
  \item Two-level obstruction principle
  (\S\ref{sec:obstruction}): separates the structural
  (nonlinear self-consistency) from the geometric
  (Widom-to-Airy transfer) difficulty.
  \item Gram-coercivity reduction
  (\S\ref{sec:gram_gap}): closes the transfer from
  edge-frame stability to a quantitative spectral gap.
\end{enumerate}

%%=================================================================
\section{Integrable Representation of the Edge Operator}%
\label{sec:iiks}
%%=================================================================

\subsection{The boundary compression operator}

\begin{definition}[$\Dedge$]\label{def:dedge}
With $\AN := \frac{1}{N}\Ev^*\Ev$ and $\Kop$ the
PSWF concentration operator:
\[
  \Dedge
  := \Pedge\,(\AN - \Kop)\,\Pedge
  : \mathrm{Edge} \to \mathrm{Edge},
\]
where $\mathrm{Edge} = \Pedge V_{N,c}$,
$\Pedge = \sum_{n=(1-\delta)N}^{N-1}\ip{\cdot}{\psi_n}\psi_n$.
\end{definition}

\begin{remark}[What $\Dedge$ encodes]
$\Dedge$ is the \emph{mismatch form} between two
quadratic form geometries restricted to the edge window:
\begin{itemize}
  \item $\AN$: arithmetic sampling geometry (discrete, prime points)
  \item $\Kop$: PSWF spectral geometry (continuous, $L^2$)
\end{itemize}
In the bulk, $\AN \approx \Kop$ (Theorem~A, Paper~XXI).
At the edge, $\Dedge$ is the precise measure of their
incompatibility. Theorem~B says $\Dedge \ne 0$;
this paper says $\Dedge \gg 0$ (positive definite).
\end{remark}

\subsection{Exact IIKS form of the sampling operator}

The key structural observation is that the prolate concentration
operator admits an exact rank-2 integrable representation,
in the sense of Its--Izergin--Korepin--Slavnov~\cite{ItsIzergin1990}.
This is not an asymptotic statement but an algebraic identity.

\begin{lemma}[Exact IIKS representation of $\Kop$]\label{lem:iiks}
Set
\[
  \Fvec(x)
  = \begin{pmatrix}e^{icx}\\e^{-icx}\end{pmatrix},
  \qquad
  \Gvec(y)
  = \begin{pmatrix}e^{-icy}\\e^{icy}\end{pmatrix},
  \qquad
  J = \begin{pmatrix}0&1\\-1&0\end{pmatrix}.
\]
Then the sinc kernel satisfies the exact rank-2 IIKS identity:
\begin{equation}\label{eq:iiks}
  K_c(x,y)
  = \frac{\Fvec(x)^\top J\,\Gvec(y)}{x-y},
\end{equation}
holding as an exact algebraic identity, with no scaling limit
or approximation involved.
\end{lemma}

\begin{proof}
Direct computation:
$\Fvec(x)^\top J\,\Gvec(y)
= e^{icx}e^{icy} - e^{-icx}e^{-icy}
= 2i\sin(c(x-y))$.
Dividing by $x-y$ and absorbing the factor $2i$ into the
normalisation of $J$ gives
$K_c(x,y) = \frac{\sin(c(x-y))}{\pi(x-y)}$,
which is \eqref{eq:iiks}.
\end{proof}

\begin{remark}[Why this matters for the transfer chain]
The IIKS form~\eqref{eq:iiks} is the \emph{only} entry point
into Riemann--Hilbert analysis for the prolate operator.
The Fourier representation alone
$K_c(x,y) = \frac{1}{2\pi}\int_{-c}^c e^{i\eta(x-y)}\,d\eta$
does not provide RHP structure.
The cubic Airy term $\tfrac{1}{3}u^3$ that governs the
edge scaling limit does not arise from the Fourier representation
but only after integrating the Widom edge phase; this integration
is visible at the RHP level (through the associated jump matrix
and $g$-function), not at the level of individual Fourier modes.
Therefore the IIKS form is not a technical convenience but a
\emph{structural necessity} for accessing the Airy geometry
of $D^{\mathrm{Airy}}$.
\end{remark}

\begin{remark}[Formal Airy limit at the RHP level (conditional)]
After the canonical edge scaling
$\eta = c + c^{1/3}u$, $x = 1 - c^{-2/3}s$,
the formal limit of the IIKS vectors under the
Deift--Zhou steepest descent analysis produces
the Airy IIKS vectors
$(\mathrm{Ai}(s), \mathrm{Ai}'(s))^\top$,
and the associated kernel converges to the Airy kernel.
This limit does \emph{not} arise from pointwise convergence
of $\Fvec$, $\Gvec$ themselves, but from convergence of the
full RHP solution $Y_c \to Y_{\mathrm{Ai}}$;
see the two-level obstruction principle in \S\ref{sec:obstruction}.
The rigorous derivation is the content of
Conjecture~\ref{conj:airy_rhp}.
\end{remark}

%%=================================================================
\section{Two-Level Obstruction Principle}\label{sec:obstruction}
%%=================================================================

The transfer from the IIKS form~\eqref{eq:iiks} to the
Airy kernel universality decomposes into two logically
independent components.
This decomposition is the conceptual core of the present paper.

\begin{center}
\fbox{\parbox{0.88\linewidth}{\centering\itshape
\textbf{Edge stability}
$=$
\textbf{(nonlinear phase consistency)}
$+$
\textbf{(linear kernel transfer)}}}
\end{center}

\medskip

\begin{definition}[Assumption~A: Phase consistency]\label{ass:A}
There exists an analytic function
$g_c:\mathbb{C}\setminus(-\infty,1]\to\mathbb{C}$
with $g_c(z) = \log z + O(z^{-1})$ as $z\to\infty$,
satisfying the nonlinear matching condition
\begin{equation}\label{eq:g-matching}
  2\operatorname{Re}(g_c^+(t))
  = \operatorname{Re}(\phi_c(t)),
  \quad t\in(-1,1),
\end{equation}
where $\phi_c = \phi_c[g_c]$ is the Widom edge phase
(depending implicitly on $g_c$ through the spectral density
near the band edge).

This is a nonlinear scalar Riemann--Hilbert problem.
Its solvability is currently open.
\end{definition}

\begin{remark}[Why Assumption~A is genuinely nonlinear]
In the classical orthogonal polynomial setting, the $g$-function
is uniquely determined by the equilibrium measure and
condition~\eqref{eq:g-matching} holds automatically as an identity.
For the prolate operator, the Widom phase $\phi_c$ is extrinsic
to any orthogonal polynomial structure and depends on $g_c$
implicitly via the spectral density $\rho$ near the band edge:
schematically $\phi_c(t) = \Phi(t, g_c(t), c)$.
This nonlinear coupling is the structural difficulty separating
the PSWF edge problem from its random matrix analogues.
No fixed candidate from random matrix theory satisfies
\eqref{eq:g-matching} for large $c$; the correct $g_c$ must
track the $O(c)$ growth of the Widom phase.
\end{remark}

\begin{definition}[Assumption~B: Steepest-descent geometry]\label{ass:B}
Given a compatible $g_c$ as in Assumption~A, the following
conditions hold for all sufficiently large $c$:
\begin{enumerate}[(i)]
  \item (Sign condition) There exist $\kappa>0$, $\epsilon>0$ such that
  $\operatorname{Re}(2g_c(z)-\phi_c(z)) \le -\kappa < 0$
  on lens contours $\Sigma_c^\pm \setminus D_\epsilon$.
  \item (Lipschitz geometry) The contours $\Sigma_c^\pm$ can be chosen
  uniformly Lipschitz with constants bounded independently of $c$.
  \item (No additional saddles)
  $\operatorname{Re}(2g_c-\phi_c)$ has no critical points on
  $\Sigma_c^\pm\setminus D_\epsilon$; the $A_2$ coalescing saddle
  at the band edge is the unique contributing saddle.
\end{enumerate}
Assumption~B is standard in form within the Deift--Zhou framework
but cannot be verified without explicit control of $g_c$
from Assumption~A.
\end{definition}

\begin{remark}[Logical separation of A and B]
Assumption~A is \emph{structural}: it asserts the existence of an
analytic object satisfying a nonlinear boundary condition.
Without Assumption~A, Assumption~B cannot even be formulated
since $2g_c - \phi_c$ is undefined.
Assumption~B is \emph{geometric}: given Assumption~A, it concerns
the sign and regularity of a known real-valued function.
The deep difficulty lies in Assumption~A;
Assumption~B is the expected next step once Assumption~A
is resolved.
\end{remark}

\begin{conjecture}[RHP convergence to Airy parametrix]%
\label{conj:airy_rhp}
Under Assumptions~A and~B, the $2\times 2$ matrix RHP
solution $Y_c$ associated to $\widetilde{K}_c^{(\xi)}$
(via the IIKS construction from Lemma~\ref{lem:iiks})
satisfies
\[
  \sup_{z\in K}\|Y_c(z) - Y_{\mathrm{Ai}}(z)\|
  \longrightarrow 0
  \quad (c\to\infty)
\]
for every compact $K\subset\mathbb{C}\setminus[0,+\infty)$,
uniformly in $\xi$ on compact sets,
where $Y_{\mathrm{Ai}}$ is the Tracy--Widom--Airy parametrix.
This implies
$\|\widetilde{K}_c^{(\xi)} - K^{\mathrm{Airy}}\|_{\mathfrak{S}_1}
= O(c^{-1/3})$
and hence Proposition~\ref{prop:airy_univ}.
\end{conjecture}

\begin{remark}[Relation to paper13cc]
Conjecture~\ref{conj:airy_rhp} is the content of
the main theorem of \cite{Paper13cc}, where the full
Deift--Zhou programme (bare RHP, $g$-function,
Airy parametrix, small-norm problem, kernel reconstruction)
is carried out explicitly, with Assumptions~A and~B
as the only remaining open inputs.
The present paper uses this as a black-box structural fact:
if Assumptions~A and~B hold, all spectral corollaries
(including the $c^{-1/3}$ gap scale) follow.
\end{remark}

%%=================================================================
\section{Gram-Coercivity Reduction Chain}\label{sec:gram_gap}
%%=================================================================

This section establishes the reduction:
\[
  \text{Edge-frame stability}
  \;\Longrightarrow\;
  \text{Gram coercivity}
  \;\Longrightarrow\;
  \text{Spectral gap lower bound}.
\]
The second implication is imported as a black-box result
from \cite{Paper13b}; the first is the new content of this paper.

\subsection{Gram coercivity at the edge}

\begin{definition}[Edge Gram matrix]
For transition-window trial vectors
$\{\phi_j^{(c)}\}_{j=1}^N \subset L^2([-1,1])$
adapted to the edge block $[(1-\delta)N, N)$, define
the Gram matrix
\[
  G_{N,c}
  = \bigl(\ip{\phi_i^{(c)}}{\phi_j^{(c)}}\bigr)_{i,j=1}^N.
\]
\end{definition}

\begin{proposition}[Reduction: frame stability $\Rightarrow$
Gram coercivity]\label{prop:frame_to_gram}
If $\Ev|_{\mathrm{Edge}}$ is a uniform quasi-isometry with
lower frame bound $A > 0$, then there exists $\alpha(c) > 0$
(depending on $A$ and the edge window width) such that
\[
  \lambda_{\min}(G_{N,c}) \ge \alpha(c) \asymp c^{-1/2}.
\]
\end{proposition}

\begin{proof}[Proof sketch]
The quasi-isometry condition controls the evaluation map
$\Ev: \mathrm{Edge} \to \mathbb{C}^N$ uniformly from below.
The Gram matrix $G_{N,c}$ is the Gram matrix of
$\{\Ev^*\delta_{p_j}\}$ restricted to the edge subspace;
the lower frame bound directly controls $\lambda_{\min}(G_{N,c})$
via the spectral representation of the frame operator.
The scale $c^{-1/2}$ comes from the Landau--Widom counting
law for the edge window.
\end{proof}

\begin{theorem}[Coercive gap transfer, after \cite{Paper13b}]%
\label{thm:gram_gap}
Assume $\lambda_{\min}(G_{N,c}) \ge \alpha(c) \asymp c^{-1/2}$
for the edge trial space.
Then there exist $C>0$ and $c_0 > 0$ such that for all
$c \ge c_0$ and all transition-window indices $n$,
\begin{equation}\label{eq:gap_lower}
  |\lambda_n^{(c)} - \lambda_{n+1}^{(c)}|
  \ge C\, c^{-1/2}.
\end{equation}
In particular, the transition spectrum is non-collapsed.
\end{theorem}

\begin{proof}
This is Theorem~A of \cite{Paper13b}, applied to the
edge trial spaces.
The Gram lower bound prevents any collapse of adjacent
compressed eigenvalues via the min-max principle;
the compressed eigenvalues control the PSWF eigenvalues
through the coercive transfer of \cite{Paper13b}~\S3.
\end{proof}

\begin{remark}[This paper is not re-proving \cite{Paper13b}]
Theorem~\ref{thm:gram_gap} uses \cite{Paper13b} as a black box.
The new content of XXII is the first implication in the chain:
showing that uniform edge-frame stability (the open question
from \texttt{HEBELSTELLE.md}) implies the Gram coercivity input
that \cite{Paper13b} requires.
The gap lower bound is a \emph{consequence}, not a goal.
\end{remark}

\begin{remark}[Conditional Airy-scale upgrade]
Under Conjecture~\ref{conj:airy_rhp}, the gap scale upgrades
from $c^{-1/2}$ (unconditional, Theorem~\ref{thm:gram_gap})
to $c^{-1/3}$ (conditional, Corollary~\ref{cor:airy_gap}).
The $c^{-1/3}$ scale is the Airy spacing of the transition
window and is the sharp scale expected from random matrix
universality; see also Open Problem~\ref{prob:airy_gap_compute}.
\end{remark}

%%=================================================================
\section{Airy Scaling Limit and Main Theorem}\label{sec:main}
%%=================================================================

\subsection{Airy universality}

\begin{proposition}[Airy kernel at the PSWF edge,
Paper~XVIII and Conjecture~\ref{conj:airy_rhp}]%
\label{prop:airy_univ}
As $c\to\infty$, $N = \lfloor\alpha c\rfloor$,
$\alpha < 2/\pi$:
\[
  K_N(N+s, N+t)
  \to K^{\mathrm{Airy}}(s,t)
  = \frac{\mathrm{Ai}(s)\mathrm{Ai}'(t)
        - \mathrm{Ai}'(s)\mathrm{Ai}(t)}{s-t},
\]
uniformly for $s,t$ in compact sets.
Here $K_N$ is the PSWF reproducing kernel at level $N$.
The convergence follows from Conjecture~\ref{conj:airy_rhp}
(via the IIKS structure of Lemma~\ref{lem:iiks}).
\end{proposition}

\subsection{$\Dedge$ as a Rand-Weyl operator}

In the Airy scaling limit, $\Dedge$ becomes:

\begin{definition}[Rand-Weyl operator]\label{def:rand_weyl}
The \emph{Rand-Weyl operator} at the PSWF edge is:
\[
  D^{\mathrm{Airy}}
  := \Pi_{\mathrm{Airy}}\,
     (A^{\mathrm{arith}} - K^{\mathrm{Airy}})
     \,\Pi_{\mathrm{Airy}},
\]
where $A^{\mathrm{arith}}$ is the limiting arithmetic sampling form
and $K^{\mathrm{Airy}}$ is the Airy kernel operator,
both compressed to the Airy edge subspace.
\end{definition}

\begin{conjecture}[Spectral gap of $D^{\mathrm{Airy}}$]%
\label{conj:airy_gap}
\[
  \inf\,\sigma(D^{\mathrm{Airy}}) \ge c_2^{\mathrm{Airy}} > 0.
\]
The constant $c_2^{\mathrm{Airy}}$ is a universal spectral constant
of the Airy kernel, independent of $\alpha$ and of the arithmetic input.
\end{conjecture}

\begin{remark}[Why the Airy kernel is the right object]
The Airy kernel is the universal edge limit of PSWF reproducing kernels
(via Proposition~\ref{prop:airy_univ} and Conjecture~\ref{conj:airy_rhp}).
The arithmetic sampling operator $A^{\mathrm{arith}}$ does
not share this universality: its structure is governed by
the distribution of primes.
The \emph{mismatch} between these two objects is $D^{\mathrm{Airy}}$,
and its uniform positivity is the geometric content of the
edge obstruction.
The two-level obstruction principle (\S\ref{sec:obstruction})
is what makes this mismatch computable: Assumption~A controls
the phase consistency that ensures $K^{\mathrm{Airy}}$ is the
correct limit, and Assumption~B controls the geometric transfer.
\end{remark}

\subsection{Main theorem}

\begin{theorem}[Edge Coercivity -- conditional on
Conjecture~\ref{conj:airy_gap}]%
\label{thm:edge_coercive}
Fix $\delta \in (0,1)$, $\alpha \in (0, 2/\pi)$.
If $\inf\,\sigma(D^{\mathrm{Airy}}) \ge c_2^{\mathrm{Airy}} > 0$,
then for $N$ sufficiently large:
\[
  \inf_{\substack{f \in \mathrm{Edge}\\ \|f\|=1}}
  \ip{\Dedge\,f}{f}
  \ge \frac{c_2^{\mathrm{Airy}}}{2} > 0.
\]
Consequently:
\begin{enumerate}[(i)]
  \item $\Ev|_{\mathrm{Edge}}$ is a uniform quasi-isometry
  on the edge spectral block.
  \item The obstruction constant
  $\tilde{c}_0 = \lmin^2 \cdot c_0^{\mathrm{III}}/4$
  is stable under small perturbations of the prime sampling measure.
  \item Theorem~B of Paper~XXI upgrades to:
  \[
    \inf_{n \in [(1-\delta)N,\,N)}
    \langle \Dedge\,\psi_n,\,\psi_n\rangle
    \ge c_2^{\mathrm{Airy}}/2 > 0.
  \]
\end{enumerate}
\end{theorem}

\begin{corollary}[Conditional Airy-scale gap]\label{cor:airy_gap}
Under Conjecture~\ref{conj:airy_rhp}, the spectral gap
in the transition window satisfies
\[
  |\lambda_n^{(c)} - \lambda_{n+1}^{(c)}|
  \ge C\,c^{-1/3}
\]
for all sufficiently large $c$ and all transition-window
indices $n \sim 2c/\pi$.
\end{corollary}

\begin{remark}[Unconditional content]
Without Conjecture~\ref{conj:airy_gap}, the following remains
unconditional (Paper~XXI):
\begin{itemize}
  \item $\Dedge \ne 0$: at least one positive eigenvalue.
  \item The witness $\psi_{n^*}$ satisfies
  $\ip{\Dedge\,\psi_{n^*}}{\psi_{n^*}} \ge \tilde{c}_0 > 0$.
\end{itemize}
Theorem~\ref{thm:gram_gap} adds, unconditionally,
the gap lower bound $|\lambda_n^{(c)}-\lambda_{n+1}^{(c)}|
\ge Cc^{-1/2}$, once Gram coercivity is established.
The upgrade to $\inf\,\sigma(\Dedge) > 0$ and the
$c^{-1/3}$ gap scale are the new conditional content of XXII.
\end{remark}

%%=================================================================
\section{Proof Strategies for the Open Conjectures}%
\label{sec:strategy}
%%=================================================================

\subsection{Approach~1: Uniform BW-Doubling over the Edge Block}

Extend the bandwidth-doubling obstruction (Paper~III)
from a pointwise statement ($\exists n^*$) to a uniform one:
\[
  \inf_{n \in [(1-\delta)N,\,N)}
  \Eout(|\psi_n|^2) \ge c_3(\delta,c) > 0.
\]
If provable, Proposition~\ref{prop:frame_to_gram} and
Theorem~\ref{thm:gram_gap} immediately give
$\ip{\Dedge\psi_n}{\psi_n} \ge \lmin^2 c_3 > 0$
for all $n$ in the edge block.

\subsection{Approach~2: RKHS Stability of $\Ev$ on Edge Modes}

Show directly that the Paley--Wiener frame condition holds
uniformly on Edge:
\[
  A\|f\|^2 \le \|\Ev\,f\|^2 \le B\|f\|^2
  \quad \forall f \in \mathrm{Edge},
\]
with $A > 0$ uniform in $n \in [(1-\delta)N, N)$.
This is a Kadec-$\tfrac{1}{4}$-type result for prime frequencies
in the edge band.

\subsection{Approach~3: Direct Analysis of $D^{\mathrm{Airy}}$}

Analyse $D^{\mathrm{Airy}}$ via Airy function oscillation:
\begin{itemize}
  \item $\mathrm{Ai}(t)$ has no positive zeros
  $\Rightarrow$ sampling of $|\mathrm{Ai}(t)|^2$
  by prime-scaled points is non-degenerate.
  \item Estimate $\inf\,\sigma(D^{\mathrm{Airy}})$
  via quantitative oscillation bounds on Airy functions.
  \item Transfer to finite $N$ using the convergence rate
  from Proposition~\ref{prop:airy_univ}.
\end{itemize}

%%=================================================================
\section{Open Problems}\label{sec:open}
%%=================================================================

\begin{problem}\label{prob:bwdoubling_uniform}
Prove uniform BW-Doubling over the edge block:
$\inf_{n \in [(1-\delta)N,N)} \Eout(|\psi_n|^2) \ge c_3 > 0$.
\end{problem}

\begin{problem}\label{prob:kadec_edge}
Prove a Kadec-type frame stability result for
$\Ev|_{\mathrm{Edge}}$ (uniform quasi-isometry).
\end{problem}

\begin{problem}\label{prob:airy_gap_compute}
Compute or bound
$c_2^{\mathrm{Airy}} = \inf\,\sigma(D^{\mathrm{Airy}})$
explicitly.
Is $c_2^{\mathrm{Airy}}$ related to the
first zero of the Airy function $a_1 \approx -2.338$?
\end{problem}

\begin{problem}\label{prob:g_existence}
Prove Assumption~A (Definition~\ref{ass:A}): existence of a
$g_c$ satisfying the nonlinear matching condition
$2\operatorname{Re}(g_c^+(t)) = \operatorname{Re}(\phi_c(t))$
on $(-1,1)$ for all sufficiently large $c$.
Natural approach: fixed-point analysis of the operator
$\mathcal{T}_c(g) := \mathcal{H}_{(-1,1)}
[\tfrac{1}{2}\operatorname{Re}(\phi_c[g])] + \log(\cdot)$.
\end{problem}

\begin{problem}\label{prob:rh_collapses}
Under RH: show that $\sigma(D^{\mathrm{Airy}})$ is controlled
by the zero distribution of $\zeta(s)$ on the critical line,
and characterise the RH-collapse of the coercivity gap.
\end{problem}

%%=================================================================
\begin{thebibliography}{99}

\bibitem{PaperIII}
S.~Schmalnauer,
\textit{PSWF Product Spectral Tail Estimates},
preprint, 2026.

\bibitem{PaperXVIII}
S.~Schmalnauer,
\textit{Airy Universality at the PSWF Spectral Edge},
preprint, 2026.

\bibitem{PaperXXI}
S.~Schmalnauer,
\textit{Commutativity Failure in the Arithmetic--Spectral Diagram}
(Paper~XXI), preprint, 2026.

\bibitem{Paper13b}
S.~Schmalnauer,
\textit{Spectral Gap Lower Bounds for Prolate Operators
via Local Kernel Approximation, Quasimode Transfer,
and Gram Coercivity}
(paper13b), preprint, 2026.

\bibitem{Paper13cc}
S.~Schmalnauer,
\textit{Airy Universality for Prolate Spheroidal Wave Functions:
Exact IIKS Structure and RHP Reduction to a Variational Problem}
(paper13cc), preprint, 2026.

\bibitem{Slepian1978}
D.~Slepian,
\textit{Prolate spheroidal wave functions~-- V},
Bell Syst.\ Tech.\ J.\ \textbf{57} (1978).

\bibitem{Widom1964}
H.~Widom,
\textit{Asymptotic behavior of the eigenvalues II},
Arch.\ Rational Mech.\ Anal.\ \textbf{17} (1964).

\bibitem{ItsIzergin1990}
A.~R.~Its, A.~G.~Izergin, V.~E.~Korepin, N.~A.~Slavnov,
\textit{Differential equations for quantum correlation functions},
Int.\ J.\ Mod.\ Phys.\ B \textbf{4} (1990), 1003--1037.

\bibitem{CCM2025}
A.~Connes, C.~Consani, and H.~Moscovici,
arXiv:2511.22755 (2025).

\end{thebibliography}

\end{document}
